\section{Bias Mitigation Improves Multiple Downstream Training Regimes}
\label{sec:downstream}

Sections~2--4 established that activation bias is:
\begin{itemize}
    \item statistically significant in activations,
    \item responsible for the majority of extreme activation outliers,
    \item a dominant contributor to second-moment inflation,
    \item and a structured component inside the spectral spike.
\end{itemize}

We now evaluate whether mitigating bias improves training across
multiple regimes.
The goal is not to solve a single isolated problem,
but to demonstrate that bias control has broad consequences
for large language model training.

All experiments use a $1$B-parameter Transformer trained for $100$B tokens.

\subsection{Low-Bit FP4 Training}

We first evaluate FP4 blockwise quantization.
The methods compared are:

\begin{enumerate}
    \item Raw FP4 (no conditioning)
    \item Mean-only removal
    \item Mean + spike whitening (top $1\%$)
    \item Metis (SVD-based spike suppression)
    \item FP16 baseline
\end{enumerate}

\paragraph{Training stability.}

Raw FP4 diverges early due to dynamic-range inflation.
Mean removal alone stabilizes training.
Mean + whitening further improves convergence smoothness.

Figure~\ref{fig:fp4_curves} shows training loss over $100$B tokens.

\begin{figure}[t]
\centering
\includegraphics[width=0.9\linewidth]{figures/fp4_training_curves.png}
\caption{FP4 training curves.
Mean removal stabilizes training relative to raw FP4.
Mean + whitening approaches SVD-based Metis performance.}
\label{fig:fp4_curves}
\end{figure}

\paragraph{Interpretation.}

The improvement is consistent with Section~2.2:
outliers are predominantly bias-driven.
Removing bias reduces blockwise scaling inflation.
Whitening refines the remaining spike variance.

Importantly, the majority of stabilization comes from bias removal,
not from full SVD suppression.

\subsection{Adaptive Optimization (AdamW)}

We next evaluate bias mitigation under full precision
using AdamW-style optimization.

Compared methods:

\begin{enumerate}
    \item AdamW baseline
    \item AdamW + activation mean removal
    \item Spectra (spike-aware optimizer)
\end{enumerate}

\paragraph{Training behavior.}

Figure~\ref{fig:opt_curves} shows convergence behavior.

\begin{figure}[t]
\centering
\includegraphics[width=0.9\linewidth]{figures/optimizer_training_curves.png}
\caption{Optimizer comparison.
Mean-centered AdamW improves stability and narrows the gap to Spectra.}
\label{fig:opt_curves}
\end{figure}

\paragraph{Interpretation.}

From Section~2.4,
bias inflates the extreme tail of the second-moment estimate $v$.
Mean removal reduces this inflation,
thereby increasing effective learning rates
without altering semantic covariance structure.

The improvement is strongest in early layers and early training,
where activation bias is largest.

\subsection{Cross-Regime Consistency}

The same structural phenomenon explains improvements in both settings:

\begin{itemize}
    \item In quantization, bias inflates dynamic range.
    \item In optimization, bias inflates second-moment statistics.
\end{itemize}

The mitigation is identical:
remove the coherent mean component.

This cross-regime consistency supports the central claim:
bias is a shared structural driver of multiple training pathologies.

\subsection{Summary of Downstream Effects}

\begin{table}[t]
\centering
\caption{Summary of bias mitigation effects.}
\label{tab:downstream_summary}
\begin{tabular}{lcc}
\toprule
Setting & Raw & Mean-Removed \\
\midrule
FP4 stability & unstable & stable \\
Block dynamic range & inflated & reduced \\
Adam second moment tail & heavy & reduced \\
Effective LR allocation & suppressed & improved \\
\bottomrule
\end{tabular}
\end{table}

Across both regimes,
mean removal produces consistent stabilization effects.

Spike whitening provides additional refinement,
but bias control is the dominant intervention.